The preceding order and Gram results first give a test at
the flat weight. We use them to select a nonzero plane
component, realize its symbol, and restore the weight. This
supplies the detecting value used in
Proposition~\ref{prop:residue-plane-neighbor-test}.
\paragraph{Choosing the plane component.}
\label{par:plane-flat-test}
\hypertarget{source-plane-flat-test}{}
Let a flat source polynomial \(H\) in the full \(k\) coordinates
in the exact operator span of Section~\ref{par:plane-orbit} have
integral coefficients in the chosen edge frame after scaling.
Take \(H\) primitive there. At each letter use the radix
specified by its frame position. Fix the labels \(\rho\)
and the degrees of the plane and complement quotients using a
monomial of \(H\) with unit coefficient. In this class, write
\(H_{\alpha,\beta}\) for the coefficient at plane quotient
\(\alpha\) and complement quotient \(\beta\). Write
\((\alpha_0,\beta_0)\) for the two quotients of the selected
monomial. All the letter positions are distinct. Coordinatewise, radix \(P\), digit \(\rho\), and quotient \(q\) give
\[
 m=\rho+Pq,\qquad 0\le\rho<P.
\]

Taking \(q\) from the plane or complement quotient as appropriate therefore recovers every true exponent. Denote the assembly of a plane polynomial
\(A\) and a complement polynomial \(N\) by
\(\operatorname{Asm}_\rho(A,N)\): multiply the fixed label
monomial by the two polynomials with variables raised to
their letter radices. Their degrees are fixed by the chosen
class. The factorial identity of Section~\ref{par:plane-digits}
gives the exact formula
\begin{equation}\label{eq:plane-slice-factor}
 \begin{split}
 [H,\operatorname{Asm}_\rho(A,N)]
 &= b_\rho
     \sum_{\alpha,\beta} H_{\alpha,\beta}\,
       A_\alpha\,u_{\rm pl}(\alpha)\alpha!\,
       N_\beta\,u_{\rm no}(\beta)\beta!.
 \end{split}
\end{equation}
Here \(b_\rho\) depends on the labels and the quotient
multidegrees, and \(u_{\rm pl},u_{\rm no}\) are the products
of the factorial units in the two sets of coordinates.
The unit \(u_{\rm pl}\) is the twist \(\gamma\) in the
plane Gram estimate. Both signs have the required uniform
approximations. The formula is obtained monomial by
monomial, so it holds for the whole class.

\begin{lemma}[A nonzero plane polynomial and a complementary test]\label{lem:plane-split-polynomial}
Let \(H\) be an integral primitive flat polynomial in the
operator span of Proposition~\ref{prop:plane-operator-prefix-orders}.
Use the digits and quotient degrees of one unit-coefficient
monomial to determine the label class in \eqref{eq:plane-slice-factor}.
Then there
are a nonzero plane polynomial \(F\), an integral complementary
test \(N_0\), and a scalar \(k_0\ne0\) such that
\eqref{eq:plane-slice-value} and \eqref{eq:plane-slice-depth}
hold for every integral plane polynomial \(A\) at the selected
quotient degrees. In particular the scalar in \eqref{eq:plane-slice-depth} satisfies
\[
 v(k_0)\le V_d-n_0+D_n.
\]
The indices of the complementary types and their order losses are chosen before \(A\).
\end{lemma}

\begin{proof}
\proofstep{The projections onto the complementary factors.}
On the complement take the product projections onto one
isotypic type at each of its chain cuts. These cuts are
nested within a frame and disjoint across frames.
The products over all choices of types sum to the identity.
For \(Y=\sum_\beta H_{\alpha_0,\beta}x^\beta\), there is a
product \(E\) with \(\nu(EY)\le\nu(Y)\). Otherwise
\(\nu(\sum_E EY)>\nu(Y)\) by the ultrametric inequality,
contrary to the sum being \(Y\). In particular \(EY\ne0\).
For this \(E\), choose an exponent \(\tau\) attaining
\(\nu(EY)=q_1\), and put \(u=v(\tau!)/2\). The plane
polynomial \(P\) has coefficients
\[
 P_\alpha=\left(E\sum_\beta H_{\alpha,\beta}x^\beta\right)_\tau .
\]

The selected plane coefficient has \(v(P_{\alpha_0})=q_1-u\), and
\[
 q_1\le\nu(Y)\le v(\beta_0!)/2,\qquad
 q_3:=\nu(P)\le q_1-u+v(\alpha_0!)/2 .
\]
The second bound for \(q_1\) uses the unit coefficient.
Choose \(z\) with \(v(z)=q_3\), possible because the
coefficients have valuations in \(\frac12\mathbb Z\)
and \(w^2=p\). The normalized plane polynomial \(F=P/z\)
then has \(\nu(F)=0\). The orbit calculation gives the signed chain and link
orders for precisely this coefficient \(P\) of the projected
complement. Scalar multiplication by \(z^{-1}\) preserves them.

\proofstep{Scalar orders in the selected plane.}
Let
\[
 n_{0,{\rm no}}=\tfrac12v(\prod_i d_{{\rm no},i}!)
\]

for the complement degrees, and let \(D_{\rm no}\) be an
integer loss bound in the factorial norm for \(E\).
It is the sum of the cut projector losses, rounded up.
There is a scalar \(s\) of valuation
\[
 n_{0,{\rm no}}+D_{\rm no}-u.
\]

The complement polynomial \(s E x^\tau\) is integral
by the weighted projection bound. It has the upper torus
weight bound of the chosen type at each cut. Apply a binomial
approximation to \(u_{\rm no}^{-1}\) with integral coefficients
of degree \(\le 2(c_P-1)\) modulo \(p^{c_P}\).
This gives an integral complement test \(N_0\). In the upper
weight convention a diagonal binomial of degree \(a\) raises
the bound by at most \(a\), by applying the Laurent estimate
with the reciprocal torus variable. Thus the upper bounds of the chosen types increase by at most the permitted allowance \(4c_P\).

If one uses the exact inverse twist in place of the
approximation, Fock self-adjointness of \(E\) in
\eqref{eq:plane-slice-factor} gives exactly
\(k_0[F,\gamma A]\) with
\[
 k_0=b_\rho s\,\tau!z.
\]

Let \(V_d=v(\prod_i d_i!)\) be the full factorial valuation.
For each exponent in full multidegree let \(\kappa_{\rm fl}\)
bound its factorial deficiency from \(V_d\).
For example \(\sum_i k\lfloor\log_p d_i\rfloor\) works,
with a zero term if \(d_i=0\), by counting carries.
When \(A\) is integral the complement approximation gives
\begin{align}
 v\bigl([H,\operatorname{Asm}_\rho(A,N_0)]-k_0[F,\gamma A]\bigr)
       &\ge c_P+V_d-\kappa_{\rm fl}, \label{eq:plane-slice-value}\\
 v(k_0)&\le V_d-n_0+D_n, \label{eq:plane-slice-depth}
\end{align}
where \(n_0\) belongs to the selected plane degrees and
\[
 D_n=D_{\rm no}+
        n_0-\tfrac12v(\alpha_0!)+
        n_{0,{\rm no}}-\tfrac12v(\beta_0!).
\]
Indeed the factorial of the selected monomial has
valuation \(v(b_\rho)+v(\alpha_0!)+v(\beta_0!)\le V_d\).
Using \(v(z)=q_3\le v(\alpha_0!)/2+v(\beta_0!)/2-u\)
in the expression for \(v(k_0)\) gives
\eqref{eq:plane-slice-depth}. Assembly preserves coefficient
valuation bounds, giving \eqref{eq:plane-slice-value}
from the full factorial lower bound.

\end{proof}

\paragraph{Using a realised flat symbol.}
Apply the plane Gram estimate with conductor \(c_P\), which
is an integer here. Its chain and link losses \(12c_P\) and
\(30c_P\) also satisfy the looser allowances \(24c_P\) and
\(60c_P\) of the flat construction. In the signed-chain
hypothesis the slack \(26rc_P\) suffices for the \(13rc_P\)
in \eqref{eq:plane-gram-reserves}. Let \(A=\phi Z_*\)
and write \(D_b=D_m+D_l+3K_d\). The construction before
the last simplification in the Gram estimate gives
\[
 v([F,\gamma A])\le n_0+D_l+1+\lceil D_m+K_d/2\rceil
                 \le n_0+D_b .
\]
Here the integral conductor data have \(K_d\ge1\).

If \(c_P>D_n+D_b+\kappa_{\rm fl}\), the error in
\eqref{eq:plane-slice-value} is above the depth of the
main term. We therefore obtain
\[
 0\ \le\
 v\left(\frac{p^{\kappa_{\rm fl}}}
              {\prod_i d_i!}[H,\operatorname{Asm}_\rho(A,N_0)]\right)
 \ \le\ D_n+D_b+\kappa_{\rm fl}.
\]
The left inequality uses coefficient integrality and the
factorial deficiency termwise. No error polynomial is
needed for this application of the Gram estimate. The bounds and decomposition below apply uniformly
to all choices of the primitive source.

The detecting value must have a depth bound common to all
primitive source combinations. We choose this bound from the
weight data before choosing a detecting polynomial. Lemma~\ref{lem:lift-plane-symbol}
then lifts the resulting flat symbol to the integral polynomial lattice.
\paragraph{A depth fixed before the choice of a source.}
\label{par:plane-uniform}
\hypertarget{source-plane-return}{}
Let
\[
 L=Nk^2\log_p a.
\]
Here \(a\) denotes the lower bound for
the flat degrees on the tops, as in Section~\ref{par:plane-rank}.
By the inequality on the selected parameters \(a\ge\max(p,4N,k,2)\), and the full
degrees are at most \(2a\). There are at most \(N\) true letters
in the signed model: their positions in the cycle are distinct tops.
Each quotient degree is at most \(2a\) divided by its low radix.
These statements hold before the digits, projections or unit
monomial of the source have been chosen.

\begin{proposition}[An integral lift of the detecting symbol at uniform depth and order]\label{prop:plane-uniform-symbol-return}
Under the hypotheses on the frames and the parameter sizes, let \(X\) be the normal degree bound in the proof and put
\[
 D_X=\lceil49Nk^2\log_p a\rceil,\qquad t_F=\lfloor X\rfloor+Nc_0.
\]
 For each primitive polynomial
in the exact model span there is a detecting flat symbol
of order at most \(t_F\) at depth \(D_X\), realizable
at the selected plane with competitor conductor \(c_P\).
Both \(D_X\) and \(t_F\) are fixed by the weight and size
data before choosing the primitive polynomial.
\end{proposition}

\begin{proof}
\proofstep{The common bounds before choosing a source.}
In \eqref{eq:plane-slice-depth} take for \(D_{\rm no}\) the sum of
the integral projector exponents. At a complement cut of total
degree \(D\), the trace projector bound gives
\(\binom K2\lfloor\log_p(2D+K-1)\rfloor\), interpreted as zero
for a scalar action. The estimate also holds in the factorial
norm by self-adjointness. There are at most \(2N\) such cuts,
since the segment with \(h+r\) listed tops has \(h+r+1\) cuts.
Now
\[
 2D+K\le4Na+k\le a^3;
\]
hence \(D_{\rm no}\le6L\).
Counting carries at each true letter bounds the plane and
complement factorial deficiencies by \(3Nr\log_p a\) and
\(3NK\log_p a\), respectively. Their halves and the
projector loss thus give \(D_n\le9L\).
The full factorial loss \(\kappa_{\rm fl}\) is at most
\(Nk\log_p(2a)\le 2Nk\log_p a\le20L\).

Similarly \(e_d\le4k^2\log_p a\) and
\(K_d\le4Nk\log_p a\), since \(2\sum_i d_i+r\le a^3\).
There are at most \(N\) links and at most \(N\) chain cuts,
so \[
 D_b=D_m+D_l+3K_d\le20L.
\] We may take the integer
\[
                         D_X=\lceil49L\rceil\le50L .
\]
Choose the absolute constant \(C\) in \(D_p\ge CL\) before
taking the host sufficiently large. The conductor hypothesis gives
\[
 c_P\ge c_\star\ge8(D_p+1)>\max\{2D_X,4D_b\},
 \qquad D_X\le50L,\quad D_b\le20L .
\]

For the second reserve in \eqref{eq:plane-gram-reserves}, write
\(D=z_j-a_{j,1}\) and \(P=p^{z_{j+1}-z_j}\). Intervening letters
have distinct distances \(m\ge D+L_1+1\) from \(z_j\).
On a chain cut in segment \(j+1\) the oldest even predecessor
lies at or after \(a_{j,1}\). These two facts give
\[
 b_j\le 2a\sum_{m\ge D+L_1+1}p^{-m}
     \le 3a p^{-D-L_1-1},
 \qquad B_u\ge \frac{a}{p^D P}.
\]
Set \(x=B_u/(Yk^2)\), so \(s_u=\lfloor x\rfloor\).

At the tested letters \(B_u\ge a/Q\) and the host size bound
give \(s_u\ge100r^2 c_P\); in particular \(x\le2s_u\).
The spacing bound \(p^{L_1+1}\ge48rYk^2\) follows from the
host spacing with \(r\le Y\) and \(N\ge1\). Hence
\begin{equation}\label{eq:plane-cut-budget}
 \frac{rb_j}{P}\le \frac{3rB_u}{p^{L_1+1}}
       \le\frac{x}{16}\le\frac{s_u}{8},\qquad
 26r c_P\le\frac{s_u}{4},\qquad 128c_P\le s_u .
\end{equation}
The last two inequalities use \(r\ge2\). Adding the first two costs in \eqref{eq:plane-cut-budget} gives
\[
 rb_j/P+26r c_P\le3s_u/8
\]
on the same cut.

\proofstep{An integral symbol from the model polynomial.}
\label{par:plane-reassembly}
Use the integral model polynomial \(Z_*\) given by the Gram
estimate, and the complement test \(N_0\). Give each monomial
in the model the segment positions of its letters. A block letter
has its intervening position, \(L_{jm}\) and \(T_{jm}\) both have
the outgoing position of that shared row, and \(W_{jm}\) has
its incoming position in the next segment. Keep segment indices
on variables even when their positions in the cycle agree. Place
the complement variables only at the intervening and outgoing
positions. These renamings commute with the signed action on a
prefix and its inverse image. This follows on the variables;
the frame rows are all even and translation by \(n'\) preserves
parity. Coincidence of the label and tangent positions causes
no loss for this assertion.

\proofstep{Prefix orders in the model plane.}
For the renamed plane polynomial the prefix before \(x_u\)
has order at least
\[
                        t_u=T_u-s_u/2-60c_P .
\]
For \(u\le h\) the inverse image is the chain cut. The initial lower bound for the order is \(T_u-s_u/8\).
The two losses in the Gram construction are at most \(3s_u/8\) and
\(24c_P\), respectively. The resulting order is at least
\[
 (T_u-s_u/8)-(3s_u/8+24c_P)=T_u-s_u/2-24c_P\ge t_u .
\]

For \(u>h\) the inverse
image consists of the model link and the outgoing label letters
at earlier rows. These extra letters are vectors. Enlarging the
torus action to them multiplies the previous transform by a
commuting substitution with nonnegative powers, so the lower
bound
\[
 T_u-3s_u/8-60c_P
\]
suffices. For the prefix through an
intervening letter the chain cut at \(u+1\) gives \(t_{u+1}\).
The renamed complement polynomial has upper orders at most
\(k_u+4c_P\) and
\[
 k_{u+1}+4c_P
\]
before and through \(x_u\).

These bounds hold on the full polynomials over \(O\). To use them
segment by segment, take homogeneous components in each
segment letter on the plane side. The complement letter degrees
are already fixed. In one component all the permitted coefficients
in a particular segment, after fixing monomials in the other
segments, belong to the same integral submodule given by the
prefix inequalities. In fact taking such coefficients commutes
with the group of the specified segment. This submodule is
saturated in the finite free module of the given degrees:
its inequalities are vanishing conditions on Taylor or Laurent
coefficients over the fraction field. The quotient is torsion free,
hence free over the DVR.

For the fixed segment degrees, let \(M_j\) be the full
polynomial lattice and \(S_j\subset M_j\) the saturated
submodule cut out by its prefix bounds. Choose an integral
retraction
\[
 r_j:M_j\longrightarrow S_j,\qquad r_j|_{S_j}=\mathrm{id}_{S_j},
\]
possible since the quotient
is free. Fixing all monomials outside \(j\) lands in \(S_j\),
so the endomorphism \(r_j\) in factor \(j\) fixes the component
\(A\). These endomorphisms commute and hence
\[
                   A=(\bigotimes_j r_j)A
                           \in\bigotimes_j S_j.
\]
Applying the same map to a monomial expansion produces
a finite integral sum of products of permitted factors.
The identical argument on the complement introduces no
denominators there either.

\proofstep{Extracting factors with the true degrees.}
Write \(A_j,G_j\) for the factors of a resulting pair. Both have
the designated letter degrees. For an even intervening \(x_u\) take \(M=k_{u+1}+4c_P\),
the upper order of \(G_j\) through the letter. The plane
order before it is at least \(t_u\), and
\[
 t_u-M\ \ge\ m_*+s_u/2-64c_P\ \ge\ m_* .
\]

For an odd intervening \(x_u\) take \(M=k_u+4c_P\)
before the letter. The plane order through it is at least
\(t_{u+1}\) and
\(t_{u+1}-M\ge m_*+s_{u+1}/2-64c_P\ge m_*\).
The lower bounds follow from the respective inequalities
for \(T_u,T_{u+1}\).
The absolute lower bound before each letter is
\(B_u/2-s_u\), since \(T_u\ge B_u/2\).
These are exactly the two parity comparisons and absolute
bound used in the determinant extraction.

For clarity, consider the degree conditions in that extraction.
At an intervening position \(y\), let \(d_{\rm low}(y)\) be the
outer digit degree and \(d_A(y),d_G(y)\) the degrees of these two
factors. The class equation gives
\[
  d_{\rm low}(y)+p^{z_j-y}\bigl(d_A(y)+d_G(y)\bigr)\le\deg(y).
\]
Thus the separate order contributions can be paid out of their sum. At an incoming position
the only letter is \(W_{j-1,m}\), of degree \(b\), and the
complementary degree is zero.

Let \(R_{\rm low}\) be the original low radix and \(R_{\rm seg}\) the incoming segment radix at this letter. Substitution first multiplies the exponent by \(P_{j-1,m}\); passing to true degree then gives
\[
 b\ \longmapsto\ P_{j-1,m}b
      \ \longmapsto\ R_{\rm low}P_{j-1,m}b=R_{\rm seg}b .
\]
This is the contribution from the incoming letter alone. Other terms at the shared letter have nonnegative degree, so if \(d_{\rm coll}\) is its total collapsed plane degree and \(d_{\rm true}\) its true degree, then
\[
 R_{\rm seg}b\le R_{\rm low}d_{\rm coll}\le d_{\rm true}.
\] This argument applies
to every model monomial used, each of which has the prescribed
collapsed degrees. After collecting terms and projecting to a
segment multidegree the same upper bound holds. In particular
the upper allowance
\[
 1.01\deg(y)/p^{z_j-y}
\]
holds in both
places. The extraction proof in Section~\ref{sec:plane-lifts}
only uses this upper inequality when estimating the earlier rows.
It therefore applies also to components which have lost much of
their degree in the incoming frame. This is the one-sided form
of the realization lemma.

\proofstep{Recovering the original low coordinates.}
Multiply each product of segment factors by the original low
label after applying the segment power substitutions. The sum
is \(\operatorname{Asm}_\rho(\phi Z_*,N_0)\): the extra power
on an incoming row is exactly \(P_{j-1,m}\) and the label has
zero support off the tops. Within each fixed component the
result has the required full degrees.

In each true letter the substitutions multiply a segment degree by its radix. Summing over the contributing parts \(b\) therefore gives 
\[
 d_{\rm collapsed}=\sum_b P_b d_b,
\]
 where \(P_b,d_b\) are the radix and degree of the intervening, incoming or outgoing part. Restoring the label then gives the original degree
by the class equation. Its normal degree is also constant: at
an even top add its normal low digits to the low radix times the
complement quotient degree; at an odd top use its plane digits
and plane quotient degree. Sum these numbers over the tops.
If \(S\) is the component of order \(o\), the one-sided realization gives an integral \(P\) with
\[
 P|_{\rm chosen}\equiv w^o S\pmod{w^{o+c_P}},\qquad P|_{\rm competitor}\equiv0\pmod{w^{o+c_P}}
\]
at every competing plane. This proves the symbol membership and its realization on the stated one-sided range.

\proofstep{Keeping one order for all tested generators.}
Finally fix an order before choosing \(H\). The counted generators
have normal degree bounded by
\[
 X=N(\mathcal L a/(pk)+10a p^{-L_1}+kQ)
\]
by
Section~\ref{par:plane-rank}. The common factor has normal degree zero. For a counted generator \(A\) and a product of \(j\le Nc_0\) row operators \(D_1,\ldots,D_j\), the per-operator increment bound gives
\[
 \deg_{\rm normal}(D_1\cdots D_j A)
     \le \lfloor X\rfloor+j\le t_F,\qquad t_F=\lfloor X\rfloor+Nc_0,
\]
whenever the result is nonzero. Taking the span preserves this upper bound. The Fock form pairs only equal exponent tuples.

If \(A=\sum_{o\le t_F}A_o\) is a tested input and
\(Z=\sum_o Z_o\) is the assembled symbol, then the diagonal
Fock pairing gives
\[
                  [A,Z]=\sum_{o\le t_F}[A_o,Z_o] .
\]
The discarded terms have disjoint exponent support from \(A\).
After this deletion the ultrametric inequality selects a
summand with pairing valuation no greater than that of the sum. We have proved the flat symbol estimate
at the same \(D_X,t_F,c_P\) for every primitive polynomial in
the exact model span. \label{cor:plane-return-1}
\end{proof}


Since the evaluation on the sum has valuation at most \(D_X\),
one of its realizable summands of order \(o\le t_F\) has
valuation at most \(D_X\).

\paragraph{Returning to the restored weight.}

The flat symbol has now been lifted with uniform depth. We must compare its detecting covector with one for the original highest weight.

Write \(\lambda\) for the full weight, \(\mu\) for its
flat part and \(A_\eta=A_r(\eta)\) for the normalisation
exponent of a weight \(\eta\). In the notation of the
edge pairing, \(R^\eta_e=w^{-A_\eta}H^\eta_e\),
\(R^{\prime\eta}_e=w^{-A_\eta}(H^\eta_e)^*\), and the
normalised evaluation is \(b^\eta_e(u,T)=w^{2A_\eta}T(u)\).
\begin{established}[Transport maps and the uniform primitive-vector bound]\label{hyp:plane-full-weight-transport}
The operator construction needed for transport has the
following properties. Its maps \(\rho_\beta:V_\lambda\to
V_\mu\) send \(R^\lambda_e\) to \(R^\mu_e\). At every
other edge into the same target they lose at most
\[
 s=rNc_0+2rn'\mathrm{cpp}
\]
powers of \(w\) on the corresponding lattices. On the
target vertex lattices they map \(H^\lambda_+\) into
\(w^{2(A_\lambda-A_\mu)-s}H^\mu_+\). They send products
of a flat generator and a remaining-weight vector into
the allowed span of operator words on that generator.
Finally, for every primitive \(u\in R^\lambda_e\), some
\(\rho_\beta u\) lies outside \(w^{2\kappa+1}R^\mu_e\),
where \(\kappa\) is bounded uniformly before \(u\) is chosen.
\end{established}
\begin{proof}
Take the maps in Proposition~\ref{prop:plane-bounded-weight-operators}. Its divided-power construction at each tensor position preserves the selected edge lattice. The sum of the losses from the other positions is \(rNc_0+2rn'\mathrm{cpp}\). The integral dual estimate there finds an image of a primitive vector outside \(w^{2\kappa+1}R_e^\mu\); its bound on \(\kappa\) depends only on the weight spreads. The construction acts by words on the remaining tensor factors and hence has the required action on the span.
\end{proof}

\begin{proposition}[A flat covector transported to the full weight]\label{prop:plane-restore-covector}
Assume the lattice maps and primitive-vector bound of
Lemma~\ref{hyp:plane-full-weight-transport}. Suppose every
primitive flat vector in the prescribed operator span has a
flat covector of depth \(D_F\), endpoint order \(t_F\), and
competitor conductor \(c_P=c_\star+s\). Then every primitive
\(u\in R^\lambda_e\) in the span of the restored model family
has a full-weight covector whose depth \(D\), endpoint order \(o\) and competitor conductor \(c\) satisfy
\[
 D\le\kappa+D_F,\qquad o\le t_F+s,\qquad c=c_\star.
\]

\end{proposition}

\begin{proof}
\proofstep{Testing a primitive combination in the restored lattice.}
For a chosen primitive combination \(u\), select such a
\(\beta\) and write \(\rho_\beta u=w^a u_1\) with \(u_1\)
primitive and \(0\le a\le2\kappa\). The operator-word
span is over the fraction field, so it contains \(u_1\).
Transport a flat covector \(T\) by
\(\Lambda T=w^{2(A_\mu-A_\lambda)}\rho_\beta^* T\).
For example \(H^\lambda_e\) maps by \(\rho_\beta\)
into \(w^{A_\lambda-A_\mu}H^\mu_e\). Pulling back the
flat lattice \(w^{-A_\mu}(H^\mu_e)^*\) and inserting
the scalar \(w^{2(A_\mu-A_\lambda)}\) gives precisely
the required inclusion in \(R^{\prime\lambda}_e\).

The calculation on the competing edge lattices loses \(s\)
powers of \(w\); the stated vertex inclusion also gives
loss \(s\) at the target vertex. The normalised pairings obey
\[
       b^\lambda_e(u,\Lambda T)=w^a b^\mu_e(u_1,T).
\]
The pairing gains valuation \(a/2\le\kappa\), since
\(w^2=p\). At the vertex and the competing edges the lattice
bounds lose at most \(s\) powers of \(w\), while the chosen
edge retains integrality. Thus the transported test has
depth at most \(\kappa+D_F\), order at most \(t_F+s\)
and conductor at least \(c_P-s=c_\star\).  The model-to-clean
congruence at these depths was checked in
Section~\ref{par:plane-middle-weights}.

\end{proof}

The final step at a single edge is to place the detecting
functional in the target lattice of
Proposition~\ref{prop:residue-plane-neighbor-test}.
The dual-plane substitution accounts for the vanishing at
the other incoming edges.
\paragraph{Dual planes at the target.}
\label{par:plane-dual-edge}
\hypertarget{source-plane-edge-covector}{}
For clarity write the distinguished edge in adapted coordinates as
\(L_-=b W^k\), \(L_+=b D_p^{-1}W^k\), where
\(D_t=\operatorname{diag}(tI_r,I_{k-r})\). Its dual edge goes
from \(L_+^\vee=b^{-t}D_p W^k\) to
\(L_-^\vee=b^{-t}W^k\).

At each Frobenius position the middle
frame on this dual edge is the inverse transpose of the original
middle frame: in the original frame use \(bD_w^{-1}\), and in
the dual use
\[
 b^{-t}D_p D_w^{-1}=b^{-t}D_w,
\]
with the
indicated Frobenius action on \(b\). Distinct incoming edges at
\(L_+\) give distinct outgoing edges at \(L_+^\vee\) by double
duality. There they have distinct residue planes, since a lattice
between a lattice and its \(p^{-1}\)-multiple is determined by
its image in that quotient. For the standard target,
\(bD_p^{-1}\in\mathrm{GL}_k(W)\); it gives the Witt frame
whose inverse transpose is \(b^{-t}D_p\). Thus a single choice
of distinguished frame in the one-sided realization handles all
incoming competitors.

\begin{lemma}[A one-sided plane symbol gives an integral dual-edge test]\label{lem:plane-dual-edge-test}
For the edge from \(bW^k\) to \(bD_p^{-1}W^k\), use
\(bD_w^{-1}\) and \(b^{-t}D_w\) as the respective middle frames.
Let a detecting
symbol of order \(o\) be realized to precision \(o+c_P\)
at the distinguished dual edge and at least that order
at its competitors. If \(D_X<c_P/2\), the resulting flat covector on the original edge has order \(o'\), conductor \(c'\) and depth \(D'\) satisfying
\[
 o'\le o,\qquad c'\ge c_P,\qquad D'\le D_X.
\]

\end{lemma}

\begin{proof}
\proofstep{The dual pairing on the target plane.}
In flat polynomial coordinates use the bilinear form
\(p^{\kappa_{\rm fl}}[P,Q]/\prod_i d_i!\).
It is integral on coefficient lattices in inverse transpose
frames. Indeed the factorial deficiency gives this in the
standard coordinates; the substitution identity
\[
 [P(MX),Q(X)]=[P(X),Q(M^tX)]
\]
follows by writing the
pairing as a constant-coefficient differentiation at zero.
The identity applies to each letter, also when it carries the
inverse transpose representation. The scalar adjustments to the two edge substitutions give 
\[
 w^{A_\mu}w^{A_\mu}=w^{2A_\mu}=p^{A_\mu},
\]
 precisely
the normalized edge pairing for the flat weight \(\mu\), by
Section~\ref{par:plane-middle-weights}.

Take the integral polynomial realizing a detecting symbol
\(\sigma\) of order \(o\), and view it as a flat section in the
dual vertex frame. Divide it by \(w^o\), and pair by this form to
obtain \(T\). At the target its order is at most \(o\), by
integrality before division. In normalized coordinates on the
dual edge its polynomial after division is \(\sigma+w^{c_P}E\)
with integral \(E\); on any other dual edge at the same source
it belongs to \(w^{c_P}O[\mathbf X]\). The inverse transpose
calculation proves \(T\in R_e^{\prime\mu}\) and
\[
 T\in w^{c_P}R_{e'}^{\prime\mu}
\]
on every competing incoming
edge. For any normalized integral input, the model error pairs into
\(w^{c_P}O\), because the factorial rescaling is integral on
the paired edge lattices. Its valuation is at least
\[
 c_P/2,
\]

strictly larger than the detected valuation bounded by \(D_X\).
The valuation of the actual pairing therefore equals that
of the model detecting value.

\end{proof}

\paragraph{Operators for the bounded remaining weights.}
\label{par:plane-rho}
First retain only the top positions in the full weight \(\lambda\).
Write \(\lambda=\mu+\eta\) there, with \(\mu\) the flat weight
and \(\eta\) the remaining weight. Its spread at each top is at
most \(c_0\). We need two facts about the Schur lattices
of Section~\ref{sec:edge-contraction}.

First, for each of the weights \(\xi=\lambda,\mu\) there is a
contravariant identification \(\Psi_\xi:V_\xi^*\longrightarrow
V_\xi\) satisfying, simultaneously in every vertex frame \(g\),
\[
 p^{h_\xi} H_\xi(g^{-t})\ \subseteq \Psi_\xi\bigl(H_\xi(g)^*\bigr)\ \subseteq\ H_\xi(g^{-t}).
\]
In this display \(g\) denotes the tuple of the component frames
and \(H^*\) the integral dual. The same inclusions hold with
\(R_e^{\prime\xi},R_{e^\vee}^\xi\) in the middle and outside,
since the two normalizations multiply the unscaled lattices by
the same \(w^{-A_\xi}\).
Here \(h_\xi\) is the sum over positions of
\(\binom k2\lfloor\log_p(l+k-1)\rfloor\) at positive spreads
\(l\), with zero at spread zero.

\begin{proposition}[Uniform operators for the remaining weights]\label{prop:plane-bounded-weight-operators}
With the bounded top and off-top spreads, the maps
\(\rho_\beta\) of Proposition~\ref{prop:plane-restore-covector}
satisfy its three lattice conditions with
\(s=rNc_0+2rn'\mathrm{cpp}\). Every primitive vector
has an image outside \(w^{2\kappa+1}R_e^\mu\) for some
such map. On the restored model family in the prescribed range of parameters, the target depth \(D\), order \(o\) and conductor \(c\) satisfy
\[
 D\le D_X+\kappa\le54Nk^2\log_p a,\qquad o\le t_F+s,\qquad c=c_\star.
\]

\end{proposition}

\begin{proof}
\proofstep{Normalizing one position in the tensor.}
To verify the normalization at one position, remove the determinant
twist. Take the matrix-polynomial Fock form divided by the
squared Fock norm of the highest polynomial \(v\), the product
of the initial principal minors prescribed by the shape.
This norm is a positive integer. Pairing with \(v\) gives evaluation
at the identity: both vanish outside its weight by row and column
degrees, and both take value one on \(v\). The highest weight
space has dimension one. For a divided-power word applied to
\(v\), transpose the word in the pairing. It follows that its
pairing with any integral polynomial in the Schur lattice is
integral.

Conversely, integrality against the Schur lattice forces
\[
 Q\in O[\mathbf X]
\]
in the matrix coordinates.
Indeed substitute a generic matrix \(M\) in the highest polynomial
by the row action. The coefficients of the resulting polynomial
in the entries of \(M\) are integral elements of the Schur
lattice, by the minor expansion. Pair it with \(Q\) and transpose
the row action. The result is \(Q\) evaluated on the generic
transposed matrix, by the identity evaluation just proved.
Hence every coefficient of \(Q\) is integral, and saturation
of the standard monomial lattice gives the assertion.
The same argument proves nondegeneracy. The divided-power span of
\(v\) contains
\[
 p^{\binom k2\lfloor\log_p(l+k-1)\rfloor}
\]

times the Schur lattice by the trace argument of
Section~\ref{sec:edge-contraction}. Thus the integral dual for the
normalized form lies between the two asserted lattices. Restoring
the determinant twists on inverse transpose frames cancels their
two scalar factors. Tensoring the forms and integral dual bases
adds the exponents and gives \(\Psi_\xi\). Its inverse pairing
is a nonzero scalar times the matrix-polynomial Fock form.
Here \(l+k-1\le a^2\) for both weights at a sufficiently late
host: the flat degree is at most \(2a\), the residual spread
is at most \(c_0\), and the host size bound gives
\(2a+c_0+k\le a^2\). Consequently
\[
 \kappa=h_\lambda+h_\mu\le4L.
\]

Second, for residual sections \(s,\beta\) consider the flat
endomorphism \(U_{s,\beta}\) defined by
\[
  [\theta(U_{s,\beta}F),\theta H]_{\rm mat}
       =[\theta F\cdot s,\theta H\cdot\beta]_{\rm mat}.
\]
The matrix-polynomial form restricted to the image of the flat
model is nondegenerate, by the preceding argument. The flat
insertion \(\theta\) uses the first column or the cofactors.
The maps \(U_{s,\beta}\) are linear in either residual input.
At a single position their span is stable under commutators
with the row operators \(E_{ij}=X_i\partial_j\).
In fact transpose adjunction and the product rule move the
commutator to the two terms applying row operators to \(s\) and
to \(\beta\), with opposite signs and transposed indices.

For cofactors the induced row action on the flat coordinates is
the negative transpose action plus the scalar trace term, by
Jacobi's identity. If the induced operator is \(A+\lambda I\), its commutator on
an endomorphism \(U\) is
\[
 (A+\lambda I)U-U(A+\lambda I)=AU-UA .
\]
Thus the determinant twist introduces no extra row weight.
The residual space is generated by products of flag minors with
at most \(c_0\) columns. Each row occurs at most once in a
column. Thus the commutator weight of a generator
\(U_{s,\beta}\) has each coordinate between \(-c_0\) and \(c_0\):
it is the difference, in one of the two orders, of the two
residual row weights. The determinant twists cancel in this difference.

\proofstep{Operators from the diagonal action.}
For degree \(d\) the endomorphism representation of
\(\operatorname{Sym}^d\) in characteristic zero has one summand for each of the highest weights
\[
 \lambda^{(i)}=(i,0,\ldots,0,-i),\qquad 0\le i\le d.
\]
Indeed \(\operatorname{Sym}^d\otimes
(\operatorname{Sym}^d)^*\) is the bidegree \((d,d)\)
space decomposed by the powers of the contraction in
Section~\ref{sec:edge-transfer}; its harmonic quotients are
the irreducible incidence spaces described there. For summand
\(i\) the row word \(E_{1k}^i\) has that highest weight and
is killed under commutator by every raising operator. Its action on a single monomial gives
\[
 E_{1k}^i X_k^d
 =\frac{d!}{(d-i)!}X_1^i X_k^{\,d-i}\ne0
 \qquad(0\le i\le d),
\]
including the identity at \(i=0\). Words of length at
most \(i\) form a commutator-stable span by the product rule and
the matrix-unit commutator. They therefore contain the entire
summand.

Complete reducibility now implies that the span of the
\(U_{s,\beta}\) lies in words of length at most \(c_0\). Indeed
any summand present has its highest weight present, and the
weight bound just proved gives \(i\le c_0\).
At several positions expand the residual sections in pure
tensors and use factorization of the Fock form. Each term gives
at most \(c_0\) row operators at each top. Changes of flat frame
preserve this word space by conjugation. The word bound also holds when the flat degree is smaller than the allowed word length.

For \(\beta\in R_{e^\vee}^{\eta}\), write \(m_\beta\) for
multiplication by that residual section. Define
\[
   \Lambda_\beta T
       =p^{h_\lambda}\Psi_\lambda^{-1}
                      m_\beta\Psi_\mu T,\qquad
   T(\rho_\beta u)=p^{A_\eta}(\Lambda_\beta T)(u).
\]
The second equality defines a linear map by finite-dimensional
duality. The edge normalizations give
\[
 b_e^\mu(\rho_\beta u,T)=b_e^\lambda(u,\Lambda_\beta T).
\]
For \(T\in R_e^{\prime\mu}\), multiplication on the dual edge and
the inclusions for \(\Psi_\xi\) give
\[
 \Psi_\mu T\in R_{e^\vee}^\mu,\qquad
 m_\beta\Psi_\mu T\in R_{e^\vee}^\lambda,\qquad
 p^{h_\lambda}R_{e^\vee}^\lambda\subseteq
                              \Psi_\lambda(R_e^{\prime\lambda}).
\]
Thus \(\Lambda_\beta T\in R_e^{\prime\lambda}\).
In unscaled dual bases the two lattice factors \(w^{-A}\) cancel
the pairing factor \(w^{2A}\). The preceding adjoint identity
and this perfectness imply
\[
 \rho_\beta R_e^\lambda\subseteq R_e^\mu.
\]

On an input
\(u=\theta F\cdot s\), the definition of the two forms shows
that its flat output is a scalar multiple of
\(U_{s,\beta}F\) in polynomial coordinates. This proves
membership in the exact operator span.

\proofstep{Reservations for the remaining weights.}
Put \(s_T=rNc_0\). Multiplying any \(\beta\) from the chosen
dual edge by \(w^{s_T}\) puts it in the dual-edge source vertex
lattice. To see this, at each position the edge normalization
subtracts the top \(r\)-row weight sum. The largest possible
difference between that sum and a sum on any
\[
 r
\]
rows is at
most \(r\) times the weight spread, by the tableau weight bound.
This proves the claim in an adapted diagonal frame, including
determinant twists, and hence in every frame. Every outgoing normalized dual edge contains the vertex lattice.
Hence the cost outside the distinguished edge is at most \(s_T\).
For the target vertex, test against \(T\in(H_+^\mu)^*\). The vertex
inclusions and \(p=w^2\) give
\begin{equation}\label{eq:edge-target-adjoint}
 \begin{aligned}
 w^{s_T}\Lambda_\beta T &\in(H_+^\lambda)^*,\\
 T(\rho_\beta u)&=w^{2A_\eta-s_T}
          (w^{s_T}\Lambda_\beta T)(u)\in w^{2A_\eta-s_T}O
                      \qquad(u\in H_+^\lambda).
 \end{aligned}
\end{equation}
Testing every integral dual functional in \eqref{eq:edge-target-adjoint} yields
\(\rho_\beta H_+^\lambda\subseteq w^{2A_\eta-s_T}H_+^\mu\).

There is also a lower bound for this family of maps. The full
lattice on the dual edge is the integral span of products
\(y\beta\) from the flat and residual edge lattices. In a standard
tableau product, the column heights of the sum shape are the
union, with multiplicities, of those of the two shapes. Split
the columns accordingly and take a tensor of these factorizations
over the positions. The frame action, determinant twist and
edge powers all multiply. Thus for every full dual integral
functional \(G\), \(\Psi_\lambda G\) is an integral sum of these
products.

If
\[
 \Psi_\lambda G=\sum_i y_i\beta_i,
\]
take
\(T_i=\Psi_\mu^{-1}(p^{h_\mu}y_i)\), which is flat
dual integral. The multiplier in \(\Lambda_\beta\) then gives
\[
 p^\kappa G=\sum_i\Lambda_{\beta_i} T_i
\]
with
\(\kappa=h_\lambda+h_\mu\). For the normalized pairings,
\[
 p^\kappa b_e^\lambda(u,G)
       =\sum_i b_e^\mu(\rho_{\beta_i}u,T_i).
\]
Suppose \(\rho_\beta u\in w^{2\kappa+1}R_e^\mu\) for every
integral choice of \(\beta\). Testing the displayed sum against the
integral flat duals then gives, for every \(G\in R_e^{\prime\lambda}\),
\[
 p^\kappa b_e^\lambda(u,G)\in w^{2\kappa+1}O
 \quad\Longrightarrow\quad b_e^\lambda(u,G)\in wO
 \quad(p^\kappa=w^{2\kappa}).
\]
Perfectness gives \(u\in w R_e^\lambda\), a contradiction for a
primitive \(u\). Thus at least one residual test detects it.

\proofstep{Restoring the non-top tensor factor.}
Restore the non-top tensor factor, of weight \(\xi\), by taking
all contractions on it with \(p^{A_\xi}G\),
\(G\in R_e^{\prime\xi}\), before applying the top maps.
In a pair of normalized dual bases these contractions are
ordinary coordinate extractions. Some extraction of a primitive
tensor remains primitive. Thus this step adds no depth. The
spread off the tops is at most \(2\mathrm{cpp}\), so
\[
 w^{2rn'\mathrm{cpp}}G
\]
belongs to the target dual vertex
lattice and hence to all competing incoming dual edge lattices.

For this inclusion, the passage from an intermediate to a target
frame expands the first \(r\) directions once more by \(w^{-1}\);
a coordinate functional then loses at most the difference of two
\(r\)-row weight sums. On the target itself the contraction
carries the additional scalar \(p^{A_\xi}\). Adding the competing-edge loss from these positions to the
top contribution gives the budget
\[
 s=s_T+2rn'\mathrm{cpp},\qquad
 2A_\eta+2A_\xi-s=2(A_\lambda-A_\mu)-s .
\]
Here the right-hand exponent is the resulting target-vertex exponent.
The flat
degree is zero at every omitted position, so this contraction
retains the top word allowance.

Apply the adjoint calculation to the primitive restored model
combination. The three resulting bounds are
\begin{equation}\label{eq:restored-covector-budget}
 \begin{aligned}
 \text{endpoint order}&\le t_F+s\le\lceil t_{\rm II}\rceil,\\
 \text{depth}&\le D_X+\kappa\le54Nk^2\log_p a\le D_{\rm II},\\
 \text{competing conductor}&\ge c_P-s=c_\star,
                    \qquad c_\star=8(D_{\rm II}+1).
 \end{aligned}
\end{equation}
The last depth comparison and the endpoint bound in
\eqref{eq:restored-covector-budget} are the host choices of Section~\ref{sec:host-edge-tests},
with its absolute coefficient taken sufficiently large.
On the source side cleaning
the normal-degree components of the original flat generators
costs at most \(\lfloor X\rfloor\) at the source, by the
one-sided lift and division by the order of each component.
Multiplication by a remaining-weight section from the chosen
edge costs at most \(s\) at that vertex, by the same spread
estimate. Thus the source order also fits. With \(c'=c_P\)
one has \[
 c'\ge2(D_X+\kappa+1),\qquad c'\ge c_\star+s.
\]
Section~\ref{par:plane-middle-weights} transfers the pairing
to the clean combination and gives its competing source
conductors. This completes the single-edge estimate with the
order and depth fixed on the host and the stated quantifier over
primitive combinations. \label{cor:plane-edge-covector-1}
\end{proof}

